\section{Eleven integer variables with a short divisibility carrier}
\label{sec:short-carrier}
This section gives the degree-440 equation and two cheaper companions.
Its two essential reductions are to replace a fifth-power carrier by
a first-power carrier, and to replace a variable-dependent index modulus
by $4q^4$. The full supported-modulus interpolation theorem is needed for
the converse. Reusing only the older interpolation modulo the strong
ordinate would leave a gap.

\subsection{The fixed constants and the complete polynomial}
The only unknowns are
\[
 (j,C,z,h,r,v,w,T,c,n,s)\in\Z^{11}.
\]
Use the unchanged source-34 tuple
\[
 \beta,K_0,M_C,M_F,A_c,B_c,
 \qquad \beta=2^{1064921}-1,\quad
 A_c=2^{1113},\quad B_c=2^{666104}.
\]
The fourth port $M_F$ is the \emph{already shifted} mask. The six exact
numerals, their bit lengths, hashes, and executable finite reconstruction
are in \code{norm/coefficients}; their common identity is
\code{mixed-parity-unrounded-1064921}. These are constants, not six extra
unknowns or tunable parameters. The tuple must not be mixed with another
source's smaller coefficients.

Define the following polynomial abbreviations:
\begin{align*}
 J&=6j+1,&q&=\beta J+1,&W&=12w+4,\\
 X&=qW,&Y&=(2v+1)q^3,&a&=(X+1)Y,\\
 H&=4a+3,&\Delta&=(a+2)^2-1,\\
 F&=(K_0+W)C-z(q-1),\\
 R&=\bigl(q(q-F)-C\bigr)(q^2-1)+(M_C+qM_F)J,\\
 k&=R+1+4h q^4,&\eta&=c-Yk,\\
 D_{\mathrm{err}}&=(A_cF-q)^2+(B_cC-q)^2,\\
 \mathcal B&=k^2(q^2-D_{\mathrm{err}})-4q^2\eta^2,\\
 \ell&=XY^2,&D_\tau&=1+\ell(\ell+1)k^2,\\
 D&=3\mathcal B-4,&D_g&=1+4D D_\tau s^2,\\
 P&=D_\tau D_g,\\
 I&=2rc^2,&D_f&=1+\Delta I^2,\\
 Q&=\Delta^2 I^2,&\delta&=Q-1,\\
 V&=D_f(c-R+4IT)-c,&C_y&=QV^2-1,\\
 D_a&=\delta C_y,&K&=X(3c-2X)+2,\\
 z_c&=HT+K,&\mathcal A&=(z_c^2+3H^2)(z_c^3-19H^3),\\
 S&=srH,&Z&=S\bigl(n+(2Q-1)V\bigr)-\mathcal A,\\
 u&=S^2D_f,&b&=4S^2D_a,&e&=16S^2P.
\end{align*}
The equation is
\begin{equation}\label{eq:three-root}
 \boxed{
 \left[(Z^2-u-b-e)^2-4(ub+ue+be)\right]^2
       -64ubeZ^2=0.}
\end{equation}
It contains only fixed nonnegative exponents and integer arithmetic. Its
literal division-free evaluator is
\code{norm/scalar11\_q4\_center.json}: 126 operations and exact degree 440.
The following proof uses radicals to analyze the polynomial; no radicals
or concealed existential witnesses occur in the equation itself.

The constant recipe constructs 943 fixed tiles from the retained literal
430-action table, assigns the fixed bit weights and positional offsets,
and forms the two masks and convolution constant by finite sums. It never
simulates the machine. Its code and complete input data are included; the
check reconstructs all four native ports and separately checks $A_c,B_c$.
The six uncompressed numeral payloads total 615,883 bytes. Thus a short
126-operation evaluator does not mean that the full coefficient
description has only 126 symbols.

\subsection{The algebraic norm and the exceptional coordinates}
Let
\[
 \mathcal F_8(Z;u,b,e)=
 [(Z^2-u-b-e)^2-4(ub+ue+be)]^2-64ubeZ^2.
\]
In a free multiquadratic algebra this is the product of
\[
 Z-\epsilon_1\sqrt u-\epsilon_2\sqrt b-\epsilon_3\sqrt e
 \quad(\epsilon_i\in\{-1,1\}).
\]
This identity holds even when radicands coincide or are negative or zero.
An independent exact checker multiplies all eight factors and compares
all 35 collected coefficients.

The fixed affine front gives $q$ even, $q\equiv2\pmod3$, $|X|\ge64$,
$|Y|\ge4096$, $H$ odd and nonzero, and $\Delta>1$.
Also $D_\tau\ge1$ because $\ell(\ell+1)\ge0$ for every integer $\ell$.
The polynomial
\[
 \pi(x)=(x^2+3)(x^3-19)
\]
has no rational root. It has a root modulo every positive integer, by
the retained elementary prime-power and CRT proof in
\code{norm/source/PRIME\_TWELVE.md}. Since
\[
 \mathcal A=H^5\pi(T+K/H),
\]
we have $\mathcal A\ne0$ for every integer assignment.

If $S=0$, every scaled radicand in~\eqref{eq:three-root} is zero and
the polynomial is $\mathcal A^8\ne0$. Thus $r,s\ne0$ at a zero.
This does \emph{not} yet imply $c\ne0$. The order of recovery matters.
We first prove $D_a>0$ in both possible cases. If $c=0$, then $I=Q=0$
and $D_a=(-1)(-1)=1$. If $c\ne0$, then $|I|\ge2c^2$,
$Q>1$, and $D_f>|c|+1$. As $V=D_f b_0-c$ for an integer $b_0$,
either $b_0=0$ and $|V|=|c|$, or
$|V|\ge D_f-|c|>1$. Thus $C_y,\delta,D_a>0$.

In a vanishing conjugate factor, $P=D_\tau D_g$ is therefore the sole
possibly negative radicand. If it were negative its nonzero imaginary
contribution could not cancel the two real roots and real center.
Since $D_g$ is odd, it follows that $D_g>0$. If $\mathcal B\le1$, then
$D\le-1$ and $D_g\le-3$, a contradiction. Hence $\mathcal B\ge2$.
For $c=0$, however,
\[
 \mathcal B=k^2\{q^2(1-4Y^2)-D_{\mathrm{err}}\}\le0.
\]
Consequently $c\ne0$, and only now do we invoke the strong and auxiliary
Pell equations with $I\ne0$ and $Q>1$.

\subsection{Square recovery and auxiliary integrality}
The normalized vanishing factor is a rational signed sum of
\[
 \sqrt{D_f},\qquad 2\sqrt{D_a},\qquad4\sqrt P.
\]
All three radicands are positive odd integers. In the squarefree
decomposition of their roots, the integer square-root coefficients are
odd. Within a nontrivial squareclass the weights have distinct 2-adic
orders $0,1,2$; the term with least order cannot cancel. Linear independence
of distinct squarefree radicals therefore forces all three roots to be
integers. Furthermore
\[
 \gcd(D_\tau,D_g)=\gcd(D_\tau,1+4D D_\tau s^2)=1,
\]
so $D_\tau=\tau^2$ and $D_g=g^2$ separately.

The square root of $D_a$ initially gives a rational $y$ satisfying
\begin{equation}\label{eq:adjacent-norm}
 QV^2-(Q-1)y^2=1.
\end{equation}
We recall the relevant integrality argument explicitly. In lowest terms
$y=a_1/b_1$, one has $b_1^2\mid Q-1$. For $|V|>1$ the usual
adjacent-coefficient Pell transformation is
\[
 V'=(2Q-1)V-2(Q-1)y,\qquad
 y'=(2Q-1)y-2QV.
\]
After choosing signs, it satisfies $|V'|<|V|$, keeps $V'$ integral, and
preserves the reduced denominator of $y$, since $2Q-1\equiv1\pmod{b_1}$.
The norm excludes $V'=0$. Descent ends at $|V|=1$, where $|y|=1$.
Thus $b_1=1$. All normalized roots in the factor are now integers, and
the factor equation yields $S\mid\mathcal A$.

\subsection{The smaller index modulus}
The retained signed rank proof uses a first-Pell congruence modulo
$|XY|>3R+1$. We may instead use
\[
 M=4q^4>3R+1.
\]
The inequality follows from the unchanged packing bound
$0<R<q^4-2$. Since $4\mid W$,
\[
 2XY^2=2W(2v+1)^2q^7\equiv0\pmod M.
\]
The signed first Pell parameter $P_0=|2XY^2+1|$ is therefore
$\operatorname{sign}(X)\pmod M$. Its classification gives
$|k|=2\psi_{P_0}(n_0)$ and the congruence
\[
 \epsilon\,2n_0\equiv R+1\pmod M,
 \qquad
 \epsilon=\operatorname{sign}(k)\operatorname{sign}(X)^{n_0-1}.
\]
For the positive sign $2n_0\ge R+1$; for the negative sign
$2n_0\ge M-R-1>2R$. The remaining strong/auxiliary folding argument
does not use $h$ separately. It gives a main index $p=R$ with $p>n_0$.
Hence the negative branch is impossible. Both $2n_0$ and $R+1$ lie
strictly between zero and $M$, so
\[
 2n_0=R+1.
\]
Only after this equality is proved do we reduce the same recurrence
modulo $|XY|$. It restores $k\equiv R+1\pmod{XY}$ and the old integral
coordinate $h_{\mathrm{old}}=(k-R-1)/(XY)$. The downstream source
argument is consequently unchanged. Conversely the canonical positive
first-Pell witness has the new congruence modulo $M$, so
$h=(k-R-1)/M$ is integral. This changes neither the witness domain nor
the coefficient tuple.

\subsection{Why fifth-power divisibility suffices}
The strong and auxiliary norms, the actual residues
\[
 V\equiv-c\pmod f,\qquad V\equiv-R\pmod c,
\]
the condition $c^2\mid I$, and the centered budget strip give the
pre-projection conclusions of the retained signed folding theorem:
\[
 R=2t_0+1,\quad t_0\ge24,\quad
 |Y|>\frac{(|X|-3)^{t_0}}2-1.
\]
These conclusions precede the missing main projector; it is not being
assumed in order to prove itself. In particular, with $x=|X|$,
\[
 |H|>x(x-3)^{t_0}>x^{t_0+1}/2^{t_0},\qquad x\ge2^{64}.
\]
From $S\mid\mathcal A$ and $H\mid S$, reduction modulo $H$ gives
$H\mid K^5$. Define $K_q=X(X+2ac)-1$. The exact identity
\[
 K=HXc-2K_q
\]
and oddness of $H$ give $H\mid K_q^5$. The signed Pell projection
congruence is
\[
 2^R K_q\equiv\Pi\pmod H,
 \quad \Pi=(\sigma X-2^R)(2^R\sigma X+1),\quad
 \sigma=\operatorname{sign}(c).
\]
Thus $H\mid\Pi^5$. Each factor of $\Pi$ has absolute value smaller
than $2^{R+1}x$, and hence
\[
 |\Pi|^5<2^{20t_0+20}x^{10}<|H|.
\]
For the last inequality, use
\[
 64(t_0-9)-(21t_0+20)=43t_0-596\ge436.
\]
Therefore $\Pi=0$, its odd second factor is nonzero, and
$\sigma X=2^R$. This recovers the exact old projector and permits the
retained orientation, half-binomial recovery, and source restoration.
The invalid inference $H\mid K^5\Rightarrow H\mid K$ is never used.

\subsection{Completeness with the strong pair fixed}
Start with the canonical outer tuple from a finite refutation. It has
$\mathcal B>64$ and $D_\tau=\tau^2$. Moreover $D_\tau\equiv1\pmod3$,
so $D D_\tau$ is positive and $2\pmod3$, hence a nonsquare. A positive
Pell solution, squared once, supplies $s,g>0$ with
\[
 g^2=1+4D D_\tau s^2.
\]
Choose one positive strong pair with $4sc^2\mid I$ and
$f^2-\Delta I^2=1$, and set $r=I/(2c^2)$. Then $r$ is an integer
divisible by $2s$. Every prime of $sr$ divides $I$, although $sr$ need
not divide $I$.

The supported-modulus interpolation theorem retained from Round 14 says
that, for this \emph{fixed} strong pair, every residue class modulo any
positive modulus whose primes divide $I$ occurs as $T$ in an integral
auxiliary solution~\eqref{eq:adjacent-norm}. Its proof uses a Pell period
$\lambda=4m_0f$ and an index increment
$d=\operatorname{lcm}(\lambda,4I)$. For each $p\mid I$, the resulting
map $h\mapsto T(h)$ satisfies
\[
 v_p(T(h_1)-T(h_2))=v_p(h_1-h_2).
\]
It is therefore a permutation modulo each supported prime power, and
CRT supplies the specified modulus. The full proof, including the
2-adic case and the precise initial index, is retained in
\code{norm/source/POWERED\_SATURATION\_PARENT.md}, Section 5, and
was independently reviewed for this use.

The canonical source already has $H\mid K$; write $\kappa=K/H$.
Choose a root of $\pi(T+\kappa)$ modulo $sr$, and interpolate that
class while leaving $I,f,r$ fixed. Then
\[
 \frac{\mathcal A}{S}
   =\frac{H^4\pi(T+\kappa)}{sr}\in\Z.
\]
Finally choose
\[
 n=\mathcal A/S+f+2\delta y+4\tau g-(2Q-1)V.
\]
One factor of the norm vanishes. The order is outer tuple, guard pair,
fixed strong pair, intersective residue, auxiliary solution, and $n$.
This proves completeness without changing a strong witness after a
residue has been selected.

\subsection{Exact degree and the two cheaper companions}
For the primary equation the important degrees are
\[
\begin{array}{c|rrrrrrrrrr}
\text{register}&k&\mathcal B&D_\tau&P&I&S&D_f&Q&V&D_a\\\hline
\deg&5&20&30&82&3&8&18&30&23&106.
\end{array}
\]
Thus $\deg(u,b,e)=(34,122,98)$. Put
$B_0=Sn-\mathcal A$ and $C_0=S(2Q-1)V$, so $Z=B_0+C_0$.
The universal identity
\[
 b-C_0^2=S^2(-V^2-4Q+4)
\]
gives $\deg(Z^2-b)=96$. Writing $\omega=Z^2-b$, the inner bracket
of~\eqref{eq:three-root} is
\[
 (\omega-u-e)^2-4(ub+ue+be).
\]
Its unique highest-degree contribution is $-4be$, of degree 220.
The correction $64ubeZ^2$ has degree 376. Hence the total degree is
exactly 440, with leading form $16b_0^2e_0^2$. On the authentic
coefficient slice a specific attained coefficient is
\[
 [j^{264}w^{40}v^{52}h^{12}s^{12}r^{24}c^{32}z^4]P_{440}
   =2^{476}3^{306}\beta^{264}>0.
\]
The exponent sum is 440. The checker independently propagates leading
forms through the actual graph and verifies the cancellation identity.
An identical 136-operation graph exposes all cancellations and has
syntactic degree bound 440; the selected 126-operation graph denotes the
same polynomial but has the looser naive syntactic bound 488.

Keeping $k=R+1+hXY$ instead gives degree 464 and 125 operations. Removing
the center shear as well gives degree 488 and 121 operations. Translation
of the unrestricted coordinate $n$ preserves zero existence exactly.
The two auxiliary identical graphs of degrees 464 and 440 are audit
artifacts, not additional frontier claims.

One charged operation is saved by the exact identity
\[
 P=D_\tau+D(2D_\tau s)^2
   =D_\tau(1+4D D_\tau s^2).
\]
The short-modulus evaluator also shares $4q^2$ with the budget, computing
$4q^4=(4q^2)q^2$ with one additional multiplication.
The actual operation counts for the three selected graphs are
\[
\begin{array}{c|rrr|r}
\text{degree}&\text{multiplications}&\text{additions}&\text{subtractions}&\text{total}\\\hline
488&77&23&21&121\\
464&79&24&22&125\\
440&80&24&22&126.
\end{array}
\]
